% !TeX program = xelatex
\documentclass[12pt]{article}
\usepackage[spanish, es-nodecimaldot]{babel}

\usepackage{fontspec}
\setmainfont{Times New Roman}

\usepackage[letterpaper,top=2.4cm,bottom=2.4cm,left=2.4cm,right=2.4cm]{geometry}

\usepackage{setspace}
% Renglon cerrado
\singlespacing
\usepackage{xurl}
\usepackage[hidelinks]{hyperref}

\usepackage{enumitem}
\usepackage{booktabs}
\usepackage{float}
\usepackage{tabularx}
\usepackage{array}

\usepackage{titlesec}

\setlength{\parindent}{0pt}

\setlength{\parskip}{6pt}

\titleformat{\section}
{\normalfont\bfseries\fontsize{12}{14}\selectfont}
{}{0pt}{}
\titlespacing*{\section}{0pt}{6pt}{0pt}

\author{Jenny Sofía Morales}

\begin{document}
	
	\begin{center}
		
		{\bfseries\fontsize{12}{14}\selectfont
			KUBERNETES , CONTENEDORES, MICROSERVICIOS EN LA NUBE
		}
		
		\vspace{0.4cm}
		
		{\itshape J. S. Morales López}
		
		\vspace{0.2cm}
		
		{\itshape 7690-08-6790 Universidad Mariano Gálvez}
		
		{\itshape Seminario de tecnologías de información}
		
		{\itshape
			\href{mailto:jmorales115@miumg.edu.gt}
			{jmorales115@miumg.edu.gt}
		}
		
		{\itshape 
			\href{https://github.com/sofiaMorales3805/ForoAcademico-Seminario}
			{Repositorio en GitHub - Sofia Morales} 
		}
	\end{center}
	
	\vspace{0.5cm}
	\section*{Resumen}
	
	
	El presente artículo aborda la orquestación de servidores y contenedores, así como los conceptos fundamentales para implementar aplicaciones basadas en microservicios dentro de entornos de nube. Su objetivo es analizar la relación existente entre Docker, Kubernetes, OAuth 2.0 y la metodología de los Doce Factores. Docker permite empaquetar cada microservicio junto con sus dependencias dentro de un contenedor. Para ello, el Dockerfile contiene las instrucciones necesarias para construir la imagen, mientras que Docker Compose facilita la definición y ejecución de aplicaciones compuestas por varios servicios. Kubernetes, por su parte, permite administrar los contenedores mediante elementos como los Pods, Deployments y Services. Estos componentes facilitan la comunicación interna, el mantenimiento de réplicas y el escalado de las aplicaciones. Asimismo, Minikube permite crear un entorno local para aprender y probar el funcionamiento de Kubernetes, mientras que kubectl proporciona los comandos para administrar el clúster. En materia de seguridad, OAuth 2.0 permite controlar el acceso a los recursos mediante tokens de autorización. Como resultado del análisis, se identifica que estas tecnologías, complementadas con los principios de los Doce Factores, favorecen el desarrollo de aplicaciones escalables, disponibles, seguras y preparadas para su implementación en la nube.
	
	\noindent
	\textbf{Palabras clave:} Docker, microservicios, Kubernetes, OAuth2.0, .
	
	
	\section*{Desarrollo del tema}
	
	La orquestación de servidores y contenedores permite coordinar los
	componentes que forman una aplicación y administrar su ejecución de manera
	automatizada. Esta práctica resulta especialmente importante en aplicaciones
	basadas en microservicios, debido a que pueden estar compuestas por numerosos
	servicios que deben comunicarse, escalarse y mantenerse disponibles dentro de
	una infraestructura local o en la nube.
	
	\noindent\textbf{Microservicios y contenedores}
	
	La arquitectura de microservicios divide una aplicación en servicios pequeños
	que representan diferentes funciones del negocio. Cada microservicio puede
	desarrollarse, desplegarse y escalarse de manera independiente. También puede
	emplear un lenguaje de programación o un mecanismo de almacenamiento diferente,
	siempre que mantenga una comunicación claramente definida con los demás
	servicios.
	
	Esta separación puede facilitar la continuidad operativa porque una falla puede
	aislarse en un componente determinado. Sin embargo, los microservicios no
	garantizan por sí solos la alta disponibilidad, ya que algunos servicios pueden
	depender de otros para completar una operación. Por esta razón, se necesitan
	mecanismos de monitoreo, recuperación, balanceo y tolerancia a fallos.
	
	Docker permite empaquetar un microservicio junto con sus dependencias dentro de
	una imagen. El Dockerfile contiene las instrucciones necesarias para construir
	dicha imagen, mientras que Docker Compose permite definir y ejecutar una
	aplicación formada por varios contenedores. En un archivo Compose pueden
	configurarse los servicios, las redes, los volúmenes y otras condiciones
	necesarias para su ejecución.
	
	\noindent\textbf{Orquestación mediante Kubernetes}
	
	Cuando la cantidad de contenedores aumenta, su administración manual se vuelve
	más compleja. Kubernetes permite orquestar aplicaciones contenerizadas,
	controlando su despliegue, comunicación, escalamiento y recuperación. Los
	contenedores se ejecutan dentro de Pods, que constituyen las unidades
	desplegables más pequeñas administradas por Kubernetes.
	
	Un Deployment declara el estado deseado de una aplicación, incluyendo la
	cantidad de réplicas que deben permanecer activas. Si un Pod presenta una falla,
	Kubernetes puede crear otro para mantener el número solicitado. Los Services
	proporcionan un punto estable de comunicación entre los componentes, aunque los
	Pods sean reemplazados y cambien de dirección dentro del clúster.
	
	La herramienta kubectl permite administrar Kubernetes mediante comandos,
	mientras que Minikube facilita la creación de un clúster local para realizar
	pruebas y conocer sus componentes. En ambientes productivos, Kubernetes puede
	implementarse sobre servidores propios o mediante servicios administrados de
	proveedores de nube.
	
	\noindent\textbf{Autorización mediante OAuth 2.0}
	
	La comunicación entre microservicios también requiere controles de seguridad.
	OAuth 2.0 es un marco de autorización que permite conceder acceso limitado a
	recursos protegidos mediante tokens. Los permisos pueden delimitarse mediante
	alcances o scopes, evitando que una aplicación obtenga más acceso del
	necesario.
	
	OAuth 2.0 no debe describirse directamente como un mecanismo de autenticación.
	Cuando se necesita verificar la identidad del usuario e implementar un inicio
	de sesión, puede utilizarse OpenID Connect sobre OAuth 2.0. Su implementación
	debe proteger los tokens, limitar sus permisos y utilizar conexiones cifradas
	para reducir el riesgo de accesos no autorizados.
	
	\noindent\textbf{Implementación en la nube}
	
	\noindent\textbf{Metodología de los Doce Factores}
		
	\section*{Observaciones y comentarios}
	

	
	\section*{Conclusiones}
	
	\begin{enumerate}
		
		\item La arquitectura debe seleccionarse tomando en cuenta los requisitos
		funcionales, los atributos de calidad y las restricciones del proyecto.
		
		\item Una aplicación pequeña puede utilizar una arquitectura monolítica o
		en capas sin asumir innecesariamente la complejidad de un sistema
		distribuido.
		
		\item No se debe asumir la complejidad de una aplicación sin realizar
		previamente un análisis de sus usuarios, funciones, presupuesto y
		crecimiento esperado.
		
		\item Una arquitectura de microservicios puede ser apropiada para
		aplicaciones complejas que requieren escalabilidad e independencia entre
		sus servicios.
		
		\item La experiencia del equipo, el presupuesto y las necesidades de
		mantenimiento también deben evaluarse para seleccionar una arquitectura
		adecuada para la organización.
		
		
	\end{enumerate}
	
	\section*{Referencias}
	
	\begin{itemize}[
		leftmargin=1.27cm,
		label={},
		itemindent=-1.27cm,
		itemsep=6pt
		]
		
		https://docs.docker.com/
		
		https://www.databricks.com/es/blog/what-is-orchestration
		https://www.advsyscon.com/blog/service-orchestration-what-is/
		https://kubernetes.io/es/
		https://cloud.google.com/learn/what-is-kubernetes?hl=es
		https://www.ibm.com/mx-es/think/topics/kubernetes
		https://aws.amazon.com/es/microservices/
		https://cloud.google.com/learn/what-is-microservices-architecture?hl=es-419
		https://www.ibm.com/mx-es/think/topics/microservices
		https://oauth.net/2/
		https://auth0.com/es/intro-to-iam/what-is-oauth-2
		https://medium.com/@csw11235/understanding-the-12-factor-app-methodology-why-we-build-cloud-applications-the-way-we-do-7aed4a98fbd8
		https://medium.com/cloud-native-daily/what-are-12-factor-principles-of-cloud-native-microservices-eb25e95e5906
		
		
		\item Fowler, M. (2015, 1 de julio).
		\textit{Microservice trade-offs}.
		MartinFowler.com.
		\url{https://martinfowler.com/articles/microservice-trade-offs.html}
		
		\item Kazman, R., Echeverria, S., \& Ivers, J. (2023).
		\textit{A holistic view of architecture definition, evolution,
			and analysis} (CMU/SEI-2023-TR-004).
		Software Engineering Institute, Carnegie Mellon University.
		\url{https://doi.org/10.1184/R1/20722402}
		
		\item Microsoft. (2025, 14 de octubre).
		\textit{Architecture styles}.
		Azure Architecture Center.
		\url{https://learn.microsoft.com/en-us/azure/architecture/guide/architecture-styles/}
		
		\item Saraswathi, R. (s. f.).
		\textit{Tipos de arquitecturas de desarrollo de aplicaciones}.
		IBM.
		\url{https://www.ibm.com/es-es/think/topics/application-architecture-types}
	\end{itemize}
	
\end{document}